\section{第一章：核心判断与底层模型}

% 写作提示（正式成稿后删除本注释）：
% 本章回答「这件事在底层是怎么回事」。先用一个反直觉判断或场景开场，
% 再展开博主的 1~2 个底层模型/隐喻，每个模型标注出处 BV，
% 全章安排 1~2 个 insight 框，章末用 \key{} 收束并把问题交给下一章。

\subsection{从一个熟悉的场景开始}

在这里放一个目标读者一看就「对，我就是这样」的场景，然后给出博主的判断。

\begin{insight}
用 \term{博主原词} 表达的核心模型：一句话写清这个模型如何解释上面的场景。
\end{insight}

\subsection{模型如何运转}

展开机制：它在什么条件下成立、和「通常以为」的解释差别在哪。
关键论断后标注出处，如（BVxxxx）。整理者自己的延伸推理用
「顺着这个逻辑……」之类的语气带出，与博主原话区分开。

\key{章末金句：把这一章压缩成读者能转述给别人的一句话。}
